\chapter*{Abstract}
\addcontentsline{toc}{chapter}{Abstract}
\selectlanguage{english}

This Master's Thesis studies agentic AI for personalised conversational
appointment management in a simulated healthcare setting. Its primary
architecture, M5, distributes five LLM responsibilities: context, profiles,
compliance, personalisation and orchestration. It is compared with S1, a
single agent, and H, a programmed hybrid workflow. All variants share tools
and mandatory access, confirmation and execution safeguards.

The prototype uses synthetic profiles and schedules and a subset of FHIR R4
to book, query, reschedule and cancel appointments and configure reminders
without external delivery. Personalisation adapts responses and ranks
compatible slots without changing permissions or operation outcomes.
CitaLab supports interaction, call inspection and conversation review.

The evaluation runs 276 episodes with two models and a simulated user,
separate from the corpus used to adjust the instructions. Task success is
checked against the schedule, and conversation quality is scored by a two-LLM
jury checked against a blind human review of a sample and against the oracle. M5
shows no significant success advantage, while its cost is systematically
higher: twice as many calls and about 2.4 times the tokens of H, which achieves
equal or better quality. No unconfirmed writes or changes to other people's
appointments are observed, yet a correctly confirmed operation can still miss
the user's intent. Conclusions are limited
to the synthetic setting studied. Code, data and
results are available at \url{https://github.com/Andros5/citalab}.

\vspace{1em}
\noindent\textbf{Keywords:} agentic AI, personalisation, multi-agent systems,
healthcare appointments, FHIR, evaluation.
\selectlanguage{spanish}
